\documentclass{article}
\usepackage[margin=1in]{geometry} % margins
\usepackage[justification=centering]{caption}
%\usepackage{calrsfs}
\usepackage{amsmath}
\usepackage{amsfonts}
\usepackage{amssymb}
\usepackage{parskip}
\usepackage{graphicx}
\usepackage{enumerate}
\usepackage{enumitem}
\usepackage{amsthm}
\usepackage{float}
\usepackage{tikz-cd}
\usepackage{ytableau}
\usepackage{dsfont}

\renewcommand{\implies}{\Rightarrow}
\newcommand{\longimplies}{\Longrightarrow}
\newcommand{\N}{\mathbb{N}}
\newcommand{\Z}{\mathbb{Z}}
\newcommand{\Q}{\mathbb{Q}}
\newcommand{\R}{\mathbb{R}}
\newcommand{\C}{\mathbb{C}}
\newcommand{\D}{\mathbb{D}}
\renewcommand{\H}{\mathbb{H}}
\newcommand{\F}{\mathbb{F}}
\newcommand{\K}{\mathbb{K}}
\newcommand{\Lf}{\mathbb{L}}
\newcommand{\Sp}{\mathbb{S}}
\newcommand{\cl}[1]{\overline{#1}}
\renewcommand{\emptyset}{\varnothing}
\newcommand{\norm}[1]{\left\|#1\right\|}
\newcommand{\maxnorm}[1]{\left\|#1\right\|_{\text{max}}}
\newcommand{\opnorm}[1]{\left\|#1\right\|_{\mathrm{op}}}
\newcommand{\supnorm}[1]{\left\|#1\right\|_{\infty}}
\newcommand{\id}{\operatorname{id}}
\newcommand{\sspan}{\operatorname{span}}
\newcommand{\tr}{\operatorname{tr}}
\newcommand{\ch}{\operatorname{char}}
\newcommand{\poly}{\mathrm{P}}
\newcommand{\mtx}{\mathcal{M}}
\newcommand{\seq}{\mathbf{Seq}}
\newcommand{\vv}[1]{\mathbf{#1}}
\newcommand{\dirsum}{\oplus}
\newcommand{\rank}{\operatorname{rank}}
\newcommand{\im}{\operatorname{im}}

\newcommand{\cj}[1]{\overline{#1}}
\newcommand{\inprod}[2]{\left\langle#1,#2\right\rangle}
\renewcommand{\Re}{\operatorname{Re}}
\renewcommand{\Im}{\operatorname{Im}}

\newcommand{\res}{\operatorname{res}}

\newcommand{\floor}[1]{\left\lfloor#1\right\rfloor}

% Theorem environments
\usepackage{amsthm}

\theoremstyle{definition} % The "plain" style italicizes all body text.
	\newtheorem{thm}{Theorem}
		\numberwithin{thm}{section} % Theorem numbers are determined by section.
	\newtheorem{lemma}[thm]{Lemma}
	\newtheorem{prop}[thm]{Proposition}
	\newtheorem{cor}[thm]{Corollary}
	\newtheorem{conj}[thm]{Conjecture}

\theoremstyle{definition} % The "definition" style does not italicize body text..
    \newtheorem{defn}[thm]{Definition}
	\newtheorem{example}[thm]{Example}
    \newtheorem{exercise}[thm]{Exercise}
    \newtheorem{question}[thm]{Question}

\usepackage{mathtools}

\usetikzlibrary{decorations.markings,positioning,decorations.text}

\def\xr{5}
\def\yr{5}

%\usepackage{titlesec}
%\titleformat{\section}[hang]{\normalfont\fontsize{14}{0}\bfseries}{Problem 8.\thesection\ }{0pt}{}

\begin{document}

\section{Exam Practice}

\begin{question}[Practice 1]
Let $f : \mathbb{R} \to \mathbb{R}$ be a continuous function with compact support. Let $E \in \mathbb{R}$ and $\eta > 0$. Compute,
\[
\lim_{\eta \to 0^+} \mathrm{Im} \left[ \int_{\mathbb{R}} \frac{f(x)}{x - (E + i \eta)} \, dx \right]
\]
If $\eta > 0$ show that the function
\[
E \to \mathrm{Im} \left[ \int_{\mathbb{R}} \frac{f(x)}{x - (E + i \eta)} \, dx \right]
\]
is smooth.
\end{question}

\begin{proof}

\end{proof}

\begin{question}[Practice 2]
Show that a monotonic function $f : \mathbb{R} \to \mathbb{R}$ is Borel measurable.
\end{question}

\begin{proof}

\end{proof}

\begin{question}[Practice 3]
If $m(E) = 0$ what can you say about $L^p(E)$?
\end{question}

\begin{proof}

\end{proof}

\begin{question}[Practice 4]
For which $p$ are trigonometric polynomials dense in $L^p(-\pi, \pi)$?
\end{question}

\begin{proof}

\end{proof}

\begin{question}[Practice 5]
Let $(M,d)$ be a metric space and for $x \in M$ and a nonempty set $A$ define
\[
d(x, A) := \inf_{y \in A} d(x,y)
\]
Prove that $d(x, A)$ is a continuous function of $x$. Prove that $A$ is closed iff the following statement holds: for all $x \in M$ we have that $d(x, A) = 0$ iff $x \in A$.
\end{question}

\begin{proof}

\end{proof}

\begin{question}[Practice 6]
Let $f_n(x)$ be a sequence of polynomials and let $f(x) = \lim_{n \to \infty} f_n(x)$. Is $f$ necessarily continuous?
\end{question}

\begin{proof}

\end{proof}

\begin{question}[Practice 7]
Let $(M, d)$ be metric space. We say a sequence $x_n \in M$ is fast-Cauchy if there is a summable sequence $a_n > 0$ (i.e., $\sum a_n < \infty$) such that $d(x_n, x_{n+1}) < a_n$. Disprove or prove the following:
\begin{enumerate}[label=(\alph*)]
    \item A fast-Cauchy sequence is Cauchy
    \item Every Cauchy sequence has a fast-Cauchy subsequence
    \item A sequence is Cauchy iff every subsequence has a subsequence that is fast-Cauchy
\end{enumerate}
\end{question}

\begin{proof}

\end{proof}


\end{document}
    